% Reusable airport network layout and drawing commands.

\newcommand{\AirportNetworkCoordinates}{%
    % Main connected component
    \coordinate (BFS) at (2.4, -0.5);
    \coordinate (LHR) at (6.9, -1.4);
    \coordinate (AMS) at (5.0,  0.8);
    \coordinate (CDG) at (7.6,  1.8);
    
    \coordinate (FRA) at (4.0,  3.2);
    \coordinate (FCO) at (0.2,  0.8);
    \coordinate (MAD) at (4.4, -2.5);
    \coordinate (LIS) at (0.2, -2.5);

    % Secondary connected component
    \coordinate (KEF) at (0.2, 3.1);
    \coordinate (OSL) at (2.5, 3.4);

    % Isolated airports
    \coordinate (EDI) at (8.8, 0.8);
    \coordinate (ATH) at (8.8, -2.2);
}
\newcommand{\DrawAirportVertices}{%
    \node[airport] at (BFS) {BFS};
    \node[airport] at (LHR) {LHR};
    \node[airport] at (AMS) {AMS};
    \node[airport] at (CDG) {CDG};
    \node[airport] at (FRA) {FRA};
    \node[airport] at (FCO) {FCO};
    \node[airport] at (MAD) {MAD};
    \node[airport] at (LIS) {LIS};

    \node[airport] at (KEF) {KEF};
    \node[airport] at (OSL) {OSL};

    \node[airport] at (EDI) {EDI};
    \node[airport] at (ATH) {ATH};
}

\newcommand{\DrawAirportRoutes}[1][route]{%
    % Main connected component
    \draw[#1] (BFS) -- (LHR);
    \draw[#1] (BFS) -- (AMS);
    \draw[#1] (BFS) -- (FCO);

    \draw[#1] (LHR) -- (CDG);
    \draw[#1] (LHR) -- (MAD);

    \draw[#1] (AMS) -- (CDG);
    \draw[#1] (AMS) -- (FRA);

    \draw[#1] (CDG) -- (FRA);

    \draw[#1] (FRA) -- (FCO);

    \draw[#1] (FCO) -- (LIS);

    \draw[#1] (MAD) -- (LIS);

    % Secondary connected component
    \draw[#1] (KEF) -- (OSL);
}
\newcommand{\DrawAdjacencyList}{%
    % BFS
    \node[adjacency airport] (bfs) at (0,0) {BFS};
    \node[adjacency neighbour] (bfs-lhr) at (2,0) {LHR};
    \node[adjacency neighbour] (bfs-ams) at (4,0) {AMS};
    \node[adjacency neighbour] (bfs-fco) at (6,0) {FCO};

    \draw[adjacency arrow] (bfs) -- (bfs-lhr);
    \draw[adjacency arrow] (bfs-lhr) -- (bfs-ams);
    \draw[adjacency arrow] (bfs-ams) -- (bfs-fco);

    % LHR
    \node[adjacency airport] (lhr) at (0,-1) {LHR};
    \node[adjacency neighbour] (lhr-bfs) at (2,-1) {BFS};
    \node[adjacency neighbour] (lhr-cdg) at (4,-1) {CDG};
    \node[adjacency neighbour] (lhr-mad) at (6,-1) {MAD};

    \draw[adjacency arrow] (lhr) -- (lhr-bfs);
    \draw[adjacency arrow] (lhr-bfs) -- (lhr-cdg);
    \draw[adjacency arrow] (lhr-cdg) -- (lhr-mad);

    % EDI
    \node[adjacency airport] (edi) at (0,-2) {EDI};
    \node[adjacency empty, anchor=west] at (1.4,-2) {no routes};
}

\newcommand{\DrawDFSActivePath}{%
    \draw[route highlighted] (BFS) -- (LHR);
    \draw[route highlighted] (LHR) -- (CDG);
    \draw[route highlighted] (CDG) -- (AMS);
}

\newcommand{\DrawDFSProgressVertices}{%
    % Previously visited airports
    \node[
        airport visited,
        label={[order badge]135:1}
    ] at (BFS) {BFS};

    \node[
        airport visited,
        label={[order badge]-45:2}
    ] at (LHR) {LHR};

    \node[
        airport visited,
        label={[order badge]45:3}
    ] at (CDG) {CDG};

    % Current recursive call
    \node[
        airport active,
        label={[order badge]135:4}
    ] at (AMS) {AMS};

    % Airports not yet visited
    \node[airport] at (FRA) {FRA};
    \node[airport] at (FCO) {FCO};
    \node[airport] at (MAD) {MAD};
    \node[airport] at (LIS) {LIS};

    % Other connected component
    \node[airport] at (KEF) {KEF};
    \node[airport] at (OSL) {OSL};

    % Isolated airports
    \node[airport] at (EDI) {EDI};
    \node[airport] at (ATH) {ATH};
}

\newcommand{\DrawDFSCallStack}{%
    \node[
        figure annotation,
        font=\sffamily\small\bfseries
    ] at (0,0.6) {Active recursive calls};

    \node[stack frame] (call-bfs) at (0,0) {
        depthFirst(BFS)
    };

    \node[stack frame] (call-lhr) at (0,-1.0) {
        depthFirst(LHR)
    };

    \node[stack frame] (call-cdg) at (0,-2.0) {
        depthFirst(CDG)
    };

    \node[stack frame active] (call-ams) at (0,-3.0) {
        depthFirst(AMS)
    };

    \draw[stack arrow] (call-bfs) -- (call-lhr);
    \draw[stack arrow] (call-lhr) -- (call-cdg);
    \draw[stack arrow] (call-cdg) -- (call-ams);

    \node[
        figure annotation,
        anchor=west,
        text=blue!70!black
    ] at (0,-3.5) {current call};
}

\newcommand{\HighlightAirportCycle}{%
    \draw[cycle direction] (BFS) -- (LHR);
    \draw[cycle direction] (LHR) -- (CDG);
    \draw[cycle direction] (CDG) -- (AMS);
    \draw[cycle direction] (AMS) -- (BFS);
}

\newcommand{\DrawCycleVertices}{%
    % Airports forming the cycle
    \node[airport active] at (BFS) {BFS};
    \node[airport active] at (LHR) {LHR};
    \node[airport active] at (CDG) {CDG};
    \node[airport active] at (AMS) {AMS};

    % Remaining airports
    \node[airport] at (FRA) {FRA};
    \node[airport] at (FCO) {FCO};
    \node[airport] at (MAD) {MAD};
    \node[airport] at (LIS) {LIS};

    \node[airport] at (KEF) {KEF};
    \node[airport] at (OSL) {OSL};

    \node[airport] at (EDI) {EDI};
    \node[airport] at (ATH) {ATH};
}

\newcommand{\DrawVisitedExplanation}{%
    \node[cycle annotation] at (0,0) {
        \textbf{Without \texttt{visited[]}:}

        \medskip

        \texttt{BFS}
        $\rightarrow$
        \texttt{LHR}
        $\rightarrow$
        \texttt{CDG}
        $\rightarrow$
        \texttt{AMS}
        $\rightarrow$
        \texttt{BFS}
        $\rightarrow \cdots$

        \medskip

        The traversal can repeatedly return to airports that it has already explored.
    };

    \node[
        cycle annotation,
        draw=green!45!black,
        fill=green!6,
        text width=42mm
    ] at (0,-3.2) {
        \textbf{With \texttt{visited[]}:}

        \medskip

        Before following a route, DFS checks whether the destination has already been visited.

        \medskip

        Each airport is therefore processed at most once.
    };
}

\newcommand{\DrawConnectedComponentRoutes}{%
    % Main connected component
    \draw[component main route] (BFS) -- (LHR);
    \draw[component main route] (BFS) -- (AMS);
    \draw[component main route] (BFS) -- (FCO);

    % Retain the remaining routes from the main component here,
    % using component main route.
    \draw[component main route] (LHR) -- (CDG);
    \draw[component main route] (LHR) -- (MAD);

    \draw[component main route] (AMS) -- (CDG);
    \draw[component main route] (AMS) -- (FRA);

    \draw[component main route] (CDG) -- (FRA);

    \draw[component main route] (FRA) -- (FCO);

    \draw[component main route] (FCO) -- (LIS);

    \draw[component main route] (MAD) -- (LIS);


    % Secondary connected component
    \draw[component secondary route] (KEF) -- (OSL);
}

\newcommand{\DrawConnectedComponentVertices}{%
    % Main connected component
    \node[component main] at (BFS) {BFS};
    \node[component main] at (LHR) {LHR};
    \node[component main] at (AMS) {AMS};
    \node[component main] at (CDG) {CDG};
    \node[component main] at (FRA) {FRA};
    \node[component main] at (FCO) {FCO};
    \node[component main] at (MAD) {MAD};
    \node[component main] at (LIS) {LIS};

    % Secondary connected component
    \node[component secondary] at (KEF) {KEF};
    \node[component secondary] at (OSL) {OSL};

    % Isolated components
    \node[component isolated] at (EDI) {EDI};
    \node[component isolated] at (ATH) {ATH};
}

\newcommand{\DrawConnectedComponentLabels}{%
    \node[
        figure annotation,
        text=green!45!black,
        font=\sffamily\small\bfseries
    ] at (4.6,-3.45) {Main connected component};

    \node[
        figure annotation,
        text=orange!65!black,
        font=\sffamily\small\bfseries
    ] at (1.35,4.05) {Secondary component};

    \node[
        component isolated label
    ] at (8.8,-3.05) {isolated airports};
}

\newcommand{\DrawWeightedAirportNetwork}{%
    \DrawAirportRoutes[weighted route]

    % Illustrative route distances
    \path (BFS) -- (LHR)
        node[route weight, above, sloped] {510 km};

    \path (BFS) -- (AMS)
        node[route weight, above, sloped] {760 km};

    \path (BFS) -- (FCO)
        node[route weight, above, sloped] {1,950 km};

    \path (LHR) -- (CDG)
        node[route weight, above, sloped] {350 km};

    \path (LHR) -- (MAD)
        node[route weight, below, sloped] {1,260 km};

    \path (CDG) -- (AMS)
        node[route weight, above, sloped] {430 km};

    \path (KEF) -- (OSL)
        node[route weight, below, sloped] {1,750 km};
}

\newcommand{\DrawWeightedAirportVertices}{%
    \node[weighted airport] at (BFS) {BFS};
    \node[weighted airport] at (LHR) {LHR};
    \node[weighted airport] at (AMS) {AMS};
    \node[weighted airport] at (CDG) {CDG};
    \node[weighted airport] at (FRA) {FRA};
    \node[weighted airport] at (FCO) {FCO};
    \node[weighted airport] at (MAD) {MAD};
    \node[weighted airport] at (LIS) {LIS};

    \node[weighted airport] at (KEF) {KEF};
    \node[weighted airport] at (OSL) {OSL};

    \node[component isolated] at (EDI) {EDI};
    \node[component isolated] at (ATH) {ATH};
}

\newcommand{\DrawWeightedGraphExplanation}{%
    \node[weight annotation] at (0,0) {
        \textbf{Weighted graph}

        \medskip

        Each route stores an additional value called a
        \emph{weight}.

        \medskip

        A weight might represent:

        \smallskip

        \begin{tabular}{@{}l@{}}
            $\bullet$ distance\\
            $\bullet$ journey time\\
            $\bullet$ ticket cost
        \end{tabular}

        \medskip

        A shortest-path algorithm compares the total weight of
        alternative routes.
    };
}